\documentclass[UTF8,zihao=-5,twoside]{ctexart}

\usepackage[a4paper,top=24mm,bottom=24mm,left=24mm,right=24mm,headheight=14pt]{geometry}
\usepackage{fontspec}
\usepackage{xeCJK}
\usepackage{unicode-math}
\usepackage{xcolor}
\usepackage{setspace}
\usepackage{indentfirst}
\usepackage{titlesec}
\usepackage{fancyhdr}
\usepackage{graphicx}
\usepackage{float}
\usepackage{caption}
\usepackage{subcaption}
\usepackage{booktabs}
\usepackage{tabularx}
\usepackage{siunitx}
\usepackage{amsmath}
\usepackage{amssymb}
\usepackage{bm}
\usepackage{authblk}
\usepackage{enumitem}
\usepackage{tikz}
\usetikzlibrary{positioning}
\usepackage{hyperref}
\usepackage[nameinlink,noabbrev]{cleveref}
\usepackage[
  backend=biber,
  style=gb7714-2015,
  sorting=none
]{biblatex}

\addbibresource{references.bib}

\IfFontExistsTF{Times New Roman}
  {\setmainfont{Times New Roman}}
  {\setmainfont{TeX Gyre Termes}}
\IfFontExistsTF{SimSun}
  {\setCJKmainfont{SimSun}}
  {\setCJKmainfont{FandolSong-Regular}}
\IfFontExistsTF{SimHei}
  {\setCJKsansfont{SimHei}}
  {\setCJKsansfont{FandolHei-Regular}}
\IfFontExistsTF{XITS Math}
  {\setmathfont{XITS Math}}
  {\setmathfont{Latin Modern Math}}

\definecolor{paperblack}{RGB}{0,0,0}
\color{paperblack}
\hypersetup{
  colorlinks=true,
  linkcolor=paperblack,
  citecolor=paperblack,
  urlcolor=paperblack,
  pdfborder={0 0 0},
  pdftitle={岩土工程科研论文}
}

\setlength{\parindent}{2em}
\setlength{\parskip}{0pt}
\setstretch{1.35}
\setlist{nosep,leftmargin=2em}
\sisetup{detect-all,per-mode=symbol,separate-uncertainty=true}
\captionsetup{font=small,labelfont=bf,labelsep=quad}
\numberwithin{equation}{section}

\titleformat{\section}
  {\bfseries\zihao{4}}
  {\thesection}{0.7em}{}
\titleformat{\subsection}
  {\bfseries\zihao{-4}}
  {\thesubsection}{0.7em}{}
\titlespacing*{\section}{0pt}{12pt}{5pt}
\titlespacing*{\subsection}{0pt}{8pt}{4pt}

\pagestyle{fancy}
\fancyhf{}
\fancyhead[LE,RO]{\zihao{5} 岩土工程科研论文}
\fancyfoot[C]{\thepage}
\renewcommand{\headrulewidth}{0.4pt}

\newcommand{\placeholder}[1]{\textcolor{paperblack}{【#1】}}
\newcommand{\keywordsname}{关键词：}
\newcommand{\enkeywordsname}{Keywords:}

\title{\bfseries\zihao{2} \placeholder{论文题目}}
\author[1]{\zihao{-4} \placeholder{作者姓名}}
\author[2]{\zihao{-4} \placeholder{合作者}}
\affil[1]{\zihao{5} \placeholder{单位名称，城市 邮编}}
\affil[2]{\zihao{5} \placeholder{单位名称，城市 邮编}}
\date{}

\begin{document}

\maketitle
\thispagestyle{plain}

\begin{abstract}
\noindent
\textbf{摘要：}
\placeholder{按研究背景、目的、方法、结果和结论组织摘要。明确研究对象、关键方法、主要定量结果及工程意义，避免只描述研究过程。}

\noindent
\textbf{\keywordsname}
\placeholder{岩土工程；边坡稳定；数值模拟；机器学习；监测预警}
\end{abstract}

\vspace{6pt}
\begin{center}
  {\bfseries\zihao{3} \placeholder{English Title of the Paper}\par}
  \vspace{5pt}
  {\zihao{-4} \placeholder{Author Name}$^{1}$, \placeholder{Co-author}$^{2}$\par}
  {\zihao{5} $^{1}$\placeholder{Affiliation, City, Postal Code, China}\par}
\end{center}

\begin{abstract}
\noindent
\textbf{Abstract:}
\placeholder{State the background, objective, methods, principal quantitative findings, and engineering implications in concise academic English.}

\noindent
\textbf{\enkeywordsname}
\placeholder{geotechnical engineering; slope stability; numerical simulation; machine learning; early warning}
\end{abstract}

\section{引言}

岩土工程问题具有材料非均质、参数不确定和施工过程强相关等特征。引言宜依次说明工程或科学问题、国内外研究进展、现有方法局限、本文研究目标与创新点。经典边坡稳定分析可追溯至 Bishop 方法\cite{bishop1955}，数据驱动方法则可用于融合监测信息与机理模型\cite{example2026}。

本文主要工作包括：
\begin{enumerate}
  \item \placeholder{提出研究问题与总体框架}；
  \item \placeholder{建立理论、数值或试验方法}；
  \item \placeholder{完成验证、对比与敏感性分析}；
  \item \placeholder{总结工程适用条件和局限性}。
\end{enumerate}

\section{工程背景与研究对象}

\subsection{工程概况}
\placeholder{说明地理位置、工程规模、支护或基础形式、建设阶段和主要风险。涉及实际工程时，应隐去不宜公开的信息。}

\subsection{工程地质条件}
\begin{table}[H]
  \centering
  \caption{岩土体主要参数}
  \label{tab:soil-parameters}
  \begin{tabularx}{\textwidth}{Xccccc}
    \toprule
    岩土层 & 重度/(\si{\kilo\newton\per\cubic\metre}) &
    弹性模量/MPa & 泊松比 & 黏聚力/kPa & 内摩擦角/\si{\degree} \\
    \midrule
    \placeholder{土层1} & \placeholder{填写} & \placeholder{填写} &
    \placeholder{填写} & \placeholder{填写} & \placeholder{填写} \\
    \placeholder{土层2} & \placeholder{填写} & \placeholder{填写} &
    \placeholder{填写} & \placeholder{填写} & \placeholder{填写} \\
    \bottomrule
  \end{tabularx}
\end{table}

\section{理论与方法}

\subsection{总体框架}
\begin{figure}[H]
  \centering
  \begin{tikzpicture}[
    node distance=12mm and 13mm,
    box/.style={draw,rounded corners=2pt,minimum width=29mm,minimum height=10mm,align=center},
    arrow/.style={->,line width=0.7pt}
  ]
    \node[box] (data) {工程与监测数据};
    \node[box,right=of data] (model) {机理或数据模型};
    \node[box,right=of model] (result) {预测与评估};
    \node[box,below=of model] (decision) {工程决策与反馈};
    \draw[arrow] (data) -- (model);
    \draw[arrow] (model) -- (result);
    \draw[arrow] (result) |- (decision);
    \draw[arrow] (decision) -| (data);
  \end{tikzpicture}
  \caption{研究总体框架}
  \label{fig:framework}
\end{figure}

\subsection{控制方程或核心算法}
以一般非线性映射为例：
\begin{equation}
  \hat{\bm{y}}=f_{\bm{\theta}}(\bm{x}),
  \label{eq:model}
\end{equation}
其中，$\bm{x}$ 为输入特征，$\hat{\bm{y}}$ 为预测结果，$\bm{\theta}$ 为待识别参数。

若采用均方误差作为目标函数，则
\begin{equation}
  \mathcal{L}(\bm{\theta})
  =\frac{1}{N}\sum_{i=1}^{N}
  \left\lVert \bm{y}_i-f_{\bm{\theta}}(\bm{x}_i)\right\rVert_2^2.
  \label{eq:loss}
\end{equation}

\subsection{参数与边界条件}
\placeholder{说明数据来源、试验条件、网格或颗粒尺度、初始应力、边界、接触模型、本构模型、地下水、荷载路径及参数标定方法。}

\section{数值模型或试验方案}

\subsection{模型建立}
\placeholder{给出模型尺寸、单元或颗粒数量、边界条件、施工阶段和计算流程。网格或颗粒尺度应完成收敛性或敏感性说明。}

\subsection{试验与数据集}
\placeholder{说明样本数量、测点布置、采样频率、数据清洗、训练验证测试划分和防止数据泄漏的措施。}

\subsection{评价指标}
常用回归指标可表示为
\begin{align}
  \mathrm{RMSE} &=
  \sqrt{\frac{1}{N}\sum_{i=1}^{N}(y_i-\hat{y}_i)^2},\\
  R^2 &=1-\frac{\sum_{i=1}^{N}(y_i-\hat{y}_i)^2}
  {\sum_{i=1}^{N}(y_i-\bar{y})^2}.
\end{align}
分类或预警研究还应报告精确率、召回率、F1 值、误报率、漏报率及提前量。

\section{结果}

\subsection{主要结果}
\begin{table}[H]
  \centering
  \caption{不同方法结果对比}
  \label{tab:model-results}
  \begin{tabular}{lcccc}
    \toprule
    方法 & RMSE & $R^2$ & 计算时间/s & 备注 \\
    \midrule
    \placeholder{基准方法} & \placeholder{填写} & \placeholder{填写} & \placeholder{填写} & \placeholder{填写} \\
    \placeholder{本文方法} & \placeholder{填写} & \placeholder{填写} & \placeholder{填写} & \placeholder{填写} \\
    \bottomrule
  \end{tabular}
\end{table}

\subsection{参数敏感性}
\placeholder{报告关键参数变化对稳定性、变形、内力或预测性能的影响。建议给出物理解释与统计不确定性。}

\section{讨论}

\subsection{机理解释}
\placeholder{联系土体结构、应力路径、渗流、施工扰动或数据分布解释结果，避免将相关性直接等同于因果关系。}

\subsection{对比与验证}
\placeholder{与理论解、现场监测、室内试验、既有方法或独立工程案例比较。说明误差来源及适用范围。}

\subsection{局限性}
\placeholder{说明样本代表性、参数不确定性、模型假设、尺度效应、外推能力和现场实施限制。}

\section{工程应用}

\placeholder{将研究结果转化为设计参数、预警阈值、施工控制指标或决策流程，并说明应用前提。}

\section{结论}

\begin{enumerate}
  \item \placeholder{结论一，应包含主要定量结果}；
  \item \placeholder{结论二，应回答核心科学或工程问题}；
  \item \placeholder{结论三，应说明工程价值和适用条件}；
  \item \placeholder{后续研究方向，保持具体且可验证}。
\end{enumerate}

\section*{数据与代码可用性}
\placeholder{填写数据、代码、模型文件的公开地址、获取条件或不公开原因。}

\section*{利益冲突}
作者声明不存在可能影响本文工作的利益冲突。

\section*{致谢}
\placeholder{填写项目、基金、单位及协作人员信息。}

\printbibliography[heading=bibintoc,title={参考文献}]

\end{document}
